\chapter{时间、正则与命令行工具}
\label{ch:11-outline}

时间、日志和命令行参数都连接着外部输入。明确时区、格式与字段约定，才能把一段临时脚本整理成可重复运行的工具。


\begin{figure}[htbp]
  \centering
  \begin{tikzpicture}[
    book mindmap,
    level 1 concept/.append style={minimum size=22mm,text width=21mm}
  ]
    \node[text=white,minimum size=25mm,text width=19mm] {第 11 章}
      child[grow=163,level distance=38mm,concept color=DeepBlue] {
        node[text=white] {时间与时区}
      }
      child[grow=17,level distance=38mm,concept color=PythonYellow] {
        node {正则与日志}
      }
      child[grow=-90,level distance=28mm,concept color=PythonBlue!22!white] {
        node {命令行与定时}
      };
  \end{tikzpicture}
  \caption*{时间与时区、日志字段提取及命令行参数。}
\end{figure}


\section[时间处理]{时间处理（datetime）}

时间点和钟面上的字样是两回事。将 UTC 时间当成统一坐标，
展示时再换成本地时区。带时区的 datetime 才能明确说明一个钟面数字对应哪个时刻，
两个地区的同一瞬间可以显示不同的小时。

存储和跨系统传递时，保留偏移或统一用 UTC。
\code{timestamp()} 配合带时区对象使用，读回时也明确传入时区。
夏令时地区需要用 zoneinfo 的地区规则，固定偏移只描述一个偏移量。


\begin{lstlisting}[title={\code{time_point.py}}]
from datetime import datetime, timedelta, timezone

# 同一时刻，分别显示 UTC 与固定的东八区偏移。
utc = datetime(2026, 1, 1, tzinfo=timezone.utc)
local = utc.astimezone(timezone(timedelta(hours=8)))
print(utc.isoformat())
print(local.isoformat())
print(utc.timestamp() == local.timestamp())
\end{lstlisting}


\textbf{运行核对}\enspace 预期输出：
\begin{lstlisting}[language={},numbers=none]
2026-01-01T00:00:00+00:00
2026-01-01T08:00:00+08:00
True
\end{lstlisting}

\par


\section[正则表达式实战]{正则表达式（Re）实战}

正则像文字筛子，先给它规定日志长什么样，再把字段挑出来。
命名组让取值处能直接写字段名。原始字符串写法 \code{r"..."} 能减少反斜杠转义的负担。

正则匹配到“像 IP 的一串字”以后，地址是否合法仍交给专门的解析器。
下面保留原日志，再用 ipaddress 核对地址。规则越写越长时，先拆字段，
把每一层需要判断的事写清楚。


\begin{lstlisting}[title={\code{parse_address.py}}]
import ipaddress
import re

# 正则先找字段，合法性再交给 IP 解析器。
line = "at=2026-01-01T00:00:00Z ip=192.0.2.1"
match = re.search(r"ip=(?P<address>\S+)", line)
if match is not None:
    address = ipaddress.ip_address(match.group("address"))
    print(address)
\end{lstlisting}


\textbf{运行核对}\enspace 预期输出：
\begin{lstlisting}[language={},numbers=none]
192.0.2.1
\end{lstlisting}

\par


\section[脚本工具化]{脚本工具化（argparse）}

命令行参数像工具面板。路径、输出位置和开关放在入口解析，
业务函数拿到明确的值再处理。这样一份脚本就能处理不同文件，
也能让定时任务用相同命令复现。

argparse 负责显示帮助和缺少参数的提示；业务代码仍须检查路径与字段是否合法。
定时执行时，工作目录、解释器、时区和日志位置都写明确。
下面用一份固定日志跑通输入路径到解析结果的过程。

\subsection*{动手｜给日志脚本加一个输入入口}

将输入路径从代码里挪到参数里，再从日志行提取时间和 IP 字段。读到格式不符的行就跳过，方便用同一条命令跑不同日志。这里的正则负责找字段，时间解析交给 datetime。

\begin{lstlisting}[title={\code{scan_log.py}}]
import argparse
from datetime import datetime, timezone
from pathlib import Path
import re

# 只提取约定格式的时间字段和 IP 字段。
pattern = re.compile(r"^(\S+)\s+ip=(\S+)$")


def main():
    parser = argparse.ArgumentParser()
    parser.add_argument("--input", type=Path, required=True)
    args = parser.parse_args()

    for line in args.input.read_text(encoding="utf-8").splitlines():
        match = pattern.fullmatch(line)
        if match is None:
            continue
        stamp, address = match.groups()
        moment = datetime.fromisoformat(stamp)
        print(moment.astimezone(timezone.utc).isoformat(), address)


if __name__ == "__main__":
    main()
\end{lstlisting}

\textbf{运行核对}\enspace 先进入 \code{examples/ch11}，再运行下列命令。\par\noindent\code{python scan_log.py --input sample.log}\par 程序应输出一条 UTC 时间和地址。省略 \code{--input} 时，argparse 会显示用法并结束。

\begin{myunderstanding}
时间点、显示格式和时区分开处理。正则先提字段，字段是否合法交给对应的解析函数核对。
\end{myunderstanding}

\begin{pycharmpractice}
在运行配置的参数栏填写 --input sample.log，再检查工作目录。终端复现时用同样参数。
\end{pycharmpractice}

\begin{pitfall}
没有时区的时间戳转换会依赖运行机器的本地配置。给存储时间和转换入口明确写出时区。
\end{pitfall}

\subsection*{进阶边界演练}
时间戳要同时交代单位与时区语境。秒和毫秒混用时，结果可能跳到完全不合理的年份。正则提取完字段后，再用日期或 IP 解析器做校验，免得一段像样的文字冒充有效数据。
